\documentclass[a4paper,14pt,oneside,openany]{memoir}

%%% Поля (ГОСТ 7.32) %%%
\usepackage[left=30mm, right=15mm, top=20mm, bottom=20mm]{geometry}
\pagestyle{plain}
\parindent=1.25cm
\usepackage{indentfirst}

%%% Язык и шрифт %%%
\usepackage{fontspec}
\usepackage{polyglossia}
\setdefaultlanguage[indentfirst=true,forceheadingpunctuation=false]{russian}
\setotherlanguages{english}

\IfFontExistsTF{Times New Roman}{
  \setmainfont{Times New Roman}
  \newfontfamily\cyrillicfont{Times New Roman}[Script=Cyrillic]
}{
  \setmainfont[
    Path=./fonts/,
    Extension=.ttf,
    UprightFont=times,
    BoldFont=timesbd,
    ItalicFont=timesi,
    BoldItalicFont=timesbi
  ]{Times New Roman}
  \newfontfamily\cyrillicfont[
    Path=./fonts/,
    Extension=.ttf,
    UprightFont=times,
    BoldFont=timesbd,
    ItalicFont=timesi,
    BoldItalicFont=timesbi
  ]{Times New Roman}[Script=Cyrillic]
}

\IfFontExistsTF{Courier New}{
  \setmonofont{Courier New}
  \newfontfamily\cyrillicfonttt{Courier New}[Script=Cyrillic]
}{
  \setmonofont[Script=Cyrillic]{FreeMono}
}

%%% Стили заголовков %%%
\setsecnumdepth{subsection}
\renewcommand*{\chapterheadstart}{}
\renewcommand*{\printchaptername}{}
\renewcommand*{\chapnumfont}{\normalfont\bfseries}
\renewcommand*{\afterchapternum}{\hspace{1em}}
\renewcommand*{\printchaptertitle}{\normalfont\bfseries\centering\MakeUppercase}
\setbeforesecskip{20pt}
\setaftersecskip{20pt}
\setsecheadstyle{\raggedright\normalfont\bfseries}
\setbeforesubsecskip{20pt}
\setaftersubsecskip{20pt}
\setsubsecheadstyle{\raggedright\normalfont\bfseries}

%%% Оглавление %%%
\addto\captionsrussian{\renewcommand\contentsname{Содержание}}
\setrmarg{2.55em plus1fil}
\renewcommand{\aftertoctitle}{\afterchaptertitle \vspace{-\cftbeforechapterskip}}
\renewcommand*{\cftchapternumwidth}{1.5em}
\renewcommand*{\cftchapterfont}{\normalfont\MakeUppercase}
\renewcommand*{\cftchapterpagefont}{\normalfont}
\renewcommand*{\cftchapterdotsep}{\cftdotsep}
\renewcommand*{\cftdotsep}{1}
\renewcommand*{\cftchapterleader}{\cftdotfill{\cftchapterdotsep}}
\maxtocdepth{subsection}

%%% Переносы %%%
\tolerance 1414
\hbadness 1414
\emergencystretch 1.5em
\hfuzz 0.3pt
\vfuzz \hfuzz
\clubpenalty=10000
\widowpenalty=10000
\brokenpenalty=4991

%%% Русские буквы для перечислений (ГОСТ: без ё з й о ч ь ы ъ) %%%
\makeatletter
  \def\russian@Alph#1{\ifcase#1\or
     А\or Б\or В\or Г\or Д\or Е\or Ж\or
     И\or К\or Л\or М\or Н\or
     П\or Р\or С\or Т\or У\or Ф\or Х\or
     Ц\or Ш\or Щ\or Э\or Ю\or Я\else\xpg@ill@value{#1}{russian@Alph}\fi}
  \def\russian@alph#1{\ifcase#1\or
     а\or б\or в\or г\or д\or е\or ж\or
     и\or к\or л\or м\or н\or
     п\or р\or с\or т\or у\or ф\or х\or
     ц\or ш\or щ\or э\or ю\or я\else\xpg@ill@value{#1}{russian@alph}\fi}
\makeatother

%%% Рисунки и таблицы %%%
\usepackage{graphicx, caption, subcaption}
\graphicspath{{img/}}

%%% В черновике отсутствующие снимки показываются явной рамкой.
%%% Перед сдачей добавьте реальные PNG в img/ и проверьте отчёт.
\newcommand{\ReportImage}[2][]{%
  \IfFileExists{#2.png}{\includegraphics[#1]{#2.png}}{%
    \IfFileExists{#2.jpg}{\includegraphics[#1]{#2.jpg}}{%
      \PackageWarning{lab-report}{Missing screenshot: #2.png}%
      \fbox{\parbox[c][35mm][c]{0.75\textwidth}{\centering
        \textbf{СКРИНШОТ НЕ ДОБАВЛЕН}\\[2mm]
        \texttt{\detokenize{#2.png}}\\[2mm]
        Добавьте настоящий скриншот перед сдачей отчёта.}}%
    }%
  }%
}

\captionsetup[figure]{font=small,width=\textwidth,name=Рисунок,justification=centering}
\captionsetup[subfigure]{font=small}
\captionsetup[table]{singlelinecheck=false,font=small,width=\textwidth,justification=justified}
\captiondelim{ — }
\setkeys{Gin}{width=0.85\textwidth}
\renewcommand{\thesubfigure}{\asbuk{subfigure}}
\usepackage[section]{placeins}

%%% Сквозная нумерация %%%
\counterwithout{figure}{chapter}
\counterwithout{equation}{chapter}
\counterwithout{table}{chapter}

%%% Списки %%%
\usepackage{enumitem}
\renewcommand*{\labelitemi}{\normalfont{--}}
\makeatletter
  \AddEnumerateCounter{\asbuk}{\russian@alph}{}
\makeatother
\renewcommand{\labelenumii}{\asbuk{enumii})}
\renewcommand{\labelenumiii}{\arabic{enumiii})}
\setlist{noitemsep, leftmargin=*}
\setlist[1]{labelindent=\parindent}
\setlist[2]{leftmargin=\parindent}
\setlist[3]{leftmargin=\parindent}

%%% Листинги %%%
\usepackage{listings}
\usepackage{xcolor}
% В listings нет встроенного CSS во всех сборках TeX Live.
\lstdefinelanguage{CSS}{
  sensitive=false,
  morecomment=[s]{/*}{*/},
  morestring=[b]",
  alsoletter={-},
  morekeywords={color,background,background-color,font-size,display,
    margin,padding,border,border-radius,grid,flex,gap}
}
\lstset{
  basicstyle=\small\ttfamily,
  breaklines=true,
  frame=single,
  tabsize=2,
  showstringspaces=false,
  inputencoding=utf8,
  extendedchars=true
}

%%% Гиперссылки %%%
\usepackage{hyperref}
\hypersetup{
  colorlinks=true,
  linktoc=all,
  linktocpage=true,
  linkcolor=red,
  urlcolor=blue
}

%%% Доп. пакеты %%%
\usepackage{longtable,ltcaption}
\usepackage{multirow,makecell}
\usepackage{booktabs}
\usepackage{soulutf8}
\usepackage{icomma}
\usepackage{textcomp}
\usepackage{amsmath}
\usepackage{amssymb} % символ \blacksquare в задании CSS

%%% Переменные отчёта %%%
\input{config}

%%% Документ %%%
\begin{document}

\input{parts/title}

\newpage
\setcounter{page}{2}
\OnehalfSpacing*

\tableofcontents*

\input{parts/intro}
\input{parts/chap1}
\chapter{Описание стенда и подготовка}
\label{ch:stand}

\section{Codepen как онлайн-редактор}

\textbf{Codepen} (\texttt{codepen.io}) --- браузерный редактор фрагментов
HTML, CSS и JavaScript, созданный Крисом Койером и Тимом Сэбэтом в
2012~году. Сервис ориентирован на быстрые демонстрации UI-приёмов и
учебные задания: каждый <<pen>> --- это набор из трёх панелей
(\texttt{HTML}, \texttt{CSS}, \texttt{JS}) и панели \texttt{Preview},
которая мгновенно перерисовывается при каждом изменении кода.

В рамках лабораторной работы~12 используется только базовый функционал:
форк готового pen-а, редактирование разметки/стилей и сохранение
результата под собственным аккаунтом.

\section{Регистрация}

На странице \texttt{codepen.io} в правом верхнем углу --- кнопка
\texttt{Sign~Up}. Доступны три способа:
\begin{itemize}
  \item \texttt{Continue~with~GitHub} --- быстрая авторизация через
        GitHub-аккаунт (рекомендуется для разработчиков);
  \item \texttt{Continue~with~Google} --- через Google-аккаунт;
  \item Email~+~пароль --- классическая регистрация с подтверждением
        по почте.
\end{itemize}

Бесплатного тарифа \texttt{Free} достаточно: лимит --- неограниченное
число публичных pen-ов, до трёх частных за всё время существования
аккаунта; кодпен-эссеты (картинки, шрифты) могут быть подключены по
внешним URL.

% — Скриншот: окно регистрации Codepen —
\begin{figure}[h]
  \centering
  \ReportImage[width=0.85\textwidth]{img/05_codepen_fork_1_1}
  \caption{Интерфейс редактора Codepen: панели HTML/CSS/JS и Preview}
  \label{fig:codepen_ui}
\end{figure}

\section{Структура pen-а}

Каждый pen состоит из четырёх взаимосвязанных панелей:
\begin{itemize}
  \item \texttt{HTML} --- разметка тела документа. Codepen
        автоматически оборачивает её в \texttt{<!DOCTYPE~html>},
        \texttt{<html>}, \texttt{<head>}, \texttt{<body>}: писать
        обвязку вручную не требуется.
  \item \texttt{CSS} --- стили; по умолчанию подключаются как
        внутренний \texttt{<style>}.
  \item \texttt{JS} --- JavaScript-код, исполняемый после загрузки
        DOM (в текущей лабораторной не используется).
  \item \texttt{Preview} --- iframe с отрисованным результатом;
        обновляется автоматически или по \texttt{Ctrl+S}.
\end{itemize}

Через кнопку \texttt{Settings} к pen-у подключаются внешние ресурсы:
\begin{itemize}
  \item \texttt{HTML~\(\rightarrow\)~Stuff~for~<head>} --- произвольные
        теги в \texttt{<head>} (например, \texttt{<meta>}, \texttt{<link>}
        на шрифты);
  \item \texttt{CSS~\(\rightarrow\)~Add~External~Stylesheets} --- список
        URL внешних таблиц стилей; в стартовых pen-ах текущей
        работы здесь уже подключены Material~Design~Lite и базовая
        таблица Netology, поэтому в собственный CSS их повторять
        не нужно.
  \item \texttt{CSS~\(\rightarrow\)~CSS~Preprocessor} --- выбор препроцессора
        (\texttt{None}, \texttt{SCSS}, \texttt{LESS}, \texttt{Stylus});
        в этой работе остаётся значение \texttt{None}.
\end{itemize}

\section{Fork-workflow}

Стартовые pen-ы методички принадлежат автору \texttt{Andrei-Andreev-the-builder}
и доступны студенту только в режиме просмотра. Чтобы внести правки,
нужно создать собственную копию --- \emph{форк}:

\begin{enumerate}
  \item Перейти на стартовый pen по ссылке \texttt{clck.ru/...} из
        методички (таблица соответствий приведена в разделе~\ref{ch:practice}).
  \item Войти в аккаунт Codepen.
  \item Нажать кнопку \texttt{Fork} в правом верхнем углу страницы
        pen-а. Codepen создаст новый pen с тем же содержимым, но
        под текущим пользователем, и автоматически откроет его в
        режиме редактирования.
  \item Внести правки в панели \texttt{HTML} (для блоков~1.1, 1.2)
        и/или \texttt{CSS} (для блоков~2.1, 2.2). Preview
        обновляется автоматически.
  \item Нажать \texttt{Save} (или \texttt{Ctrl+S}) --- pen
        получит уникальный URL, который остаётся за студентом
        навсегда. Этот URL вставляется в раздел \texttt{<<Готовые
        форки в Codepen>>} главы~\ref{ch:practice}.
\end{enumerate}

% — Скриншот: кнопка Fork → Save в Codepen —
\begin{figure}[h]
  \centering
  \ReportImage[width=0.85\textwidth]{img/09_codepen_save}
  \caption{Кнопки \texttt{Fork} и \texttt{Save} в правой верхней панели Codepen}
  \label{fig:codepen_fork}
\end{figure}

\section{Параллельная локальная разработка}

Codepen удобен для коротких демо, но плохо подходит для редактирования
вне браузера или для контроля версий через git. Поэтому в каталоге
\texttt{lab12/codepen-solutions/} хранятся локальные копии каждого
из четырёх решений: четыре подкаталога (\texttt{1.1-html-tags/},
\texttt{1.2-lists-tables/}, \texttt{2.1-selectors/},
\texttt{2.2-text-styling/}), каждый с собственным \texttt{index.html}
и --- для блоков, требующих CSS, --- \texttt{styles.css}.

Все четыре изображения, требуемые методичкой (три фотографии гейзеров
Камчатки для блока~1.1 и фон цитаты для блока~2.2), скачаны из
методички Яндекс.Диска и положены в \texttt{codepen-solutions/img/};
ссылки в HTML --- относительные (\texttt{../img/<имя>.jpg}).
Это позволяет открывать локальные страницы офлайн: для рендеринга
требуется только CDN с шрифтами Google~Fonts, Material~Design~Lite
и общими таблицами стилей Netology (\texttt{netology-code.github.io}).

Запуск любого решения:
\begin{lstlisting}[language=bash]
xdg-open lab12/codepen-solutions/1.1-html-tags/index.html
xdg-open lab12/codepen-solutions/1.2-lists-tables/index.html
xdg-open lab12/codepen-solutions/2.1-selectors/index.html
xdg-open lab12/codepen-solutions/2.2-text-styling/index.html
\end{lstlisting}

После проверки локального рендера правки переносятся в
соответствующий форк Codepen.

\endinput

\chapter{Решение четырёх блоков}
\label{ch:practice}

Методичка делит работу на два блока (HTML и CSS), внутри каждого --- два
подблока. Итого четыре независимых задания, каждое --- отдельный pen на
Codepen. Соответствие подблоков и pen-ов:

\begin{center}
\begin{tabular}{|c|l|l|l|}
\hline
\textbf{Подблок} & \textbf{Тема} & \textbf{Pen (slug)} & \textbf{Short URL} \\
\hline
1.1 & Теги для разметки текста и атрибуты            & \texttt{wvVrwjg} & \texttt{clck.ru/3E73mG} \\
1.2 & Списки и таблицы                                & \texttt{KKOXPBZ} & \texttt{clck.ru/3E73xM} \\
2.1 & Селекторы и свойства                            & \texttt{MWNEgBG} & \texttt{clck.ru/3E744t} \\
2.2 & Оформление текстовых блоков при помощи CSS      & \texttt{eYqGOjj} & \texttt{clck.ru/3E74AH} \\
\hline
\end{tabular}
\end{center}

\section{Блок 1.1: Теги для разметки текста}

Исходный pen содержит три секции, в каждой требуется выполнить отдельную
задачу разметки.

\subsection{Задача 1 --- статья про нейросеть}

Текст про нейросеть в первом \texttt{<section>} был представлен без
структуры (просто два абзаца текста, разделённые переводами строк).
Требовалось:
\begin{enumerate}
  \item Превратить предложение \textit{<<Что такое нейросеть>>} в
        заголовок второго уровня (\texttt{<h2>}).
  \item Выделить два следующих абзаца тегами \texttt{<p>}.
  \item Во втором абзаце слово \textit{<<научилась>>} превратить
        в гиперссылку на статью о работе нейросети Яндекса со
        стихами Егора Летова.
\end{enumerate}

Готовая разметка:
\begin{lstlisting}[language=HTML]
<section>
  <h2>Что такое нейросеть</h2>
  <p>Грубо говоря, это сложная программа, состоящая из многих
     программ поменьше, ...</p>
  <p>Недавно, например, нейросеть
     <a href="https://www.sostav.ru/publication/nejroset-yandeksa
     -napisala-tekst-dlya-muzykalnogo-alboma-22860.html">научилась</a>
     писать песни в стиле Егора Летова. ...</p>
</section>
\end{lstlisting}

\subsection{Задача 2 --- разметка <<Дорога>>}

В исходном pen-е разметка второго \texttt{<section>} содержала
четыре ошибки: закрывающий \texttt{</h2>} без открывающего тега для
слова \textit{<<Дорога>>}; \texttt{<blockquote>} с двумя
закрывающими тегами \texttt{</blockquote>} и \texttt{</p>} в
неправильном порядке. Требовалось исправить разметку так, чтобы
получилась последовательность: заголовок --- абзац --- цитата ---
абзац.

Готовая разметка:
\begin{lstlisting}[language=HTML]
<section>
  <h2>Дорога</h2>
  <p>Вечером второго дня я свернул в деревню Мшага ...</p>
  <blockquote>
    <p>Вдоль трассы М10 очень много мёртвых или полумёртвых
       деревень ...</p>
  </blockquote>
  <p>После того как я проехал Великий Новгород, дожди
     закончились ...</p>
</section>
\end{lstlisting}

\subsection{Задача 3 --- гейзеры Камчатки}

Третий \texttt{<section>} содержал заголовок второго уровня
\textit{<<Источники и гейзеры Камчатки>>} (уже корректно
размеченный) и три коротких рассказа --- про посёлок Малки,
Налычевскую долину и Паратунку. Требовалось:
\begin{itemize}
  \item Превратить названия гейзеров в заголовки третьего уровня
        (\texttt{<h3>}).
  \item Описание каждого гейзера обернуть в \texttt{<p>}.
  \item Перед каждым описанием вставить тег \texttt{<img>} с
        соответствующей фотографией и атрибутом \texttt{alt},
        равным названию гейзера.
\end{itemize}

Файлы изображений лежат в \texttt{lab12/codepen-solutions/img/} и
подключаются относительными ссылками:
\begin{lstlisting}[language=HTML]
<h3>Посёлок Малки</h3>
<img src="../img/malki.jpg" alt="Посёлок Малки">
<p>Малки – посёлок в 132 км от Петропавловска-Камчатского. ...</p>

<h3>Налычевская долина</h3>
<img src="../img/nalychevo.jpg" alt="Налычевская долина">
<p>...</p>

<h3>Паратунка</h3>
<img src="../img/paratunka.jpg" alt="Паратунка">
<p>...</p>
\end{lstlisting}

% — Скриншот: рендер решения блока 1.1 —
\begin{figure}[h]
  \centering
  \ReportImage[width=0.85\textwidth]{img/01_block_1_1}
  \caption{Рендер решения блока 1.1: разметка текста, ссылка, фотографии гейзеров}
  \label{fig:block_1_1}
\end{figure}

Готовый форк pen-а на Codepen:
\begin{lstlisting}
https://codepen.io/blxckweed/pen/EaNmYge
\end{lstlisting}

\section{Блок 1.2: Списки и таблицы}

Исходный pen содержит четыре секции, в каждой --- отдельная задача на
списки или таблицы.

\subsection{Задача 1 --- маркированный список}

В первой секции были перечислены пять имён через запятую:
Мартин Лютер Кинг, Роберт Оппенгеймер, Альфред Хичкок, Чарли Чаплин,
Коко Шанель. Требовалось превратить перечисление в маркированный
список \texttt{<ul>}.

\begin{lstlisting}[language=HTML]
<ul>
  <li>Мартин Лютер Кинг</li>
  <li>Роберт Оппенгеймер</li>
  <li>Альфред Хичкок</li>
  <li>Чарли Чаплин</li>
  <li>Коко Шанель</li>
</ul>
\end{lstlisting}

\subsection{Задача 2 --- упорядоченный список миллиардеров}

ТОП-5 списка Forbes требовалось расположить в порядке \emph{убывания}
состояния. Корректный порядок: Гейтс (\$75~млрд), Ортега
(\$67~млрд), Баффет (\$60{,}8~млрд), Слим (\$50~млрд), Безос
(\$45{,}2~млрд). Поскольку порядок важен, использован
\texttt{<ol>}, а не \texttt{<ul>}.

\begin{lstlisting}[language=HTML]
<ol>
  <li>Билл Гейтс ($75 млрд)</li>
  <li>Амансио Ортега ($67 млрд)</li>
  <li>Уоррен Баффет ($60,8 млрд)</li>
  <li>Карлос Слим Элу ($50 млрд)</li>
  <li>Джефф Безос ($45,2 млрд)</li>
</ol>
\end{lstlisting}

\subsection{Задача 3 --- таблица заказа аниматоров}

Дано семь строк <<персонаж/время/цена>>. Требовалось построить
таблицу с шапкой <<Персонаж~/~Время выступления~/~Цена>> и итоговой
строкой <<ИТОГО>>, занимающей две первые ячейки, с общей суммой в
третьей. Сумма проверяется арифметически:
$1500+4800+8400+5200+3300+4000+2150=29\,350$~руб.

Также требовалось использовать <<все теги, изученные на занятии>>:
\texttt{<table>}, \texttt{<thead>}, \texttt{<tbody>},
\texttt{<tfoot>}, \texttt{<tr>}, \texttt{<th>}, \texttt{<td>},
\texttt{colspan}.

\begin{lstlisting}[language=HTML]
<table>
  <thead>
    <tr>
      <th>Персонаж</th><th>Время выступления</th><th>Цена</th>
    </tr>
  </thead>
  <tbody>
    <tr><td>Белоснежка</td><td>12:00</td><td>1 500 руб.</td></tr>
    <tr><td>Лунтик</td><td>12:30</td><td>4 800 руб.</td></tr>
    <tr><td>Маша и Медведь</td><td>13:15</td><td>8 400 руб.</td></tr>
    <tr><td>Бэтмен</td><td>13:45</td><td>5 200 руб.</td></tr>
    <tr><td>Джек Воробей</td><td>14:30</td><td>3 300 руб.</td></tr>
    <tr><td>Аватар</td><td>14:50</td><td>4 000 руб.</td></tr>
    <tr><td>Энгри Бердс</td><td>15:20</td><td>2 150 руб.</td></tr>
  </tbody>
  <tfoot>
    <tr>
      <td colspan="2">ИТОГО</td>
      <td>29 350 руб.</td>
    </tr>
  </tfoot>
</table>
\end{lstlisting}

\subsection{Задача 4 --- починка таблицы звёзд эстрады}

В исходной таблице были три типа ошибок:
\begin{enumerate}
  \item Лишние пустые ячейки во второй строке шапки --- из-за
        \texttt{rowspan="2"} в первой строке колонки <<Возраст>>
        и <<Знак зодиака>> уже занимают вторую строку, и
        пустые \texttt{<td>} там избыточны.
  \item Лишний пустой \texttt{<td>} в строке Челентано.
  \item Атрибут \texttt{rowspan="2"} у ячейки <<Близнецы>> в
        строке Азнавура без соответствующего сокращения числа
        ячеек в строке Маккартни --- из-за этого следующая
        строка получает лишний \texttt{<td>}, который смещает
        остальные данные.
  \item Лишний пустой \texttt{<td>} в строке Мадонны.
\end{enumerate}

Исправленная таблица --- ровно семь строк и четыре колонки
(2 строки шапки + 5 строк данных):

\begin{lstlisting}[language=HTML]
<table>
  <thead>
    <tr>
      <th colspan="2">ФИО</th>
      <th rowspan="2">Возраст</th>
      <th rowspan="2">Знак зодиака</th>
    </tr>
    <tr>
      <th>Имя</th>
      <th>Фамилия</th>
    </tr>
  </thead>
  <tbody>
    <tr><td>Майкл</td><td>Джексон</td><td>51</td><td>Дева</td></tr>
    <tr><td>Адриано</td><td>Челентано</td><td>78</td><td>Козерог</td></tr>
    <tr><td>Шарль</td><td>Азнавур</td><td>94</td><td>Близнецы</td></tr>
    <tr><td>Пол</td><td>Маккартни</td><td>74</td><td>Близнецы</td></tr>
    <tr><td colspan="2">Мадонна</td><td>58</td><td>Лев</td></tr>
  </tbody>
</table>
\end{lstlisting}

% — Скриншот: рендер решения блока 1.2 —
\begin{figure}[h]
  \centering
  \ReportImage[width=0.85\textwidth]{img/02_block_1_2}
  \caption{Рендер решения блока 1.2: два списка и две таблицы}
  \label{fig:block_1_2}
\end{figure}

\begin{figure}[h]
  \centering
  \ReportImage[width=0.85\textwidth]{img/06_codepen_fork_1_2}
  \caption{Форк второго HTML-задания в Codepen}
  \label{fig:codepen_fork_12}
\end{figure}

Готовый форк pen-а на Codepen:
\begin{lstlisting}
https://codepen.io/blxckweed/pen/YPpVKpm
\end{lstlisting}

\section{Блок 2.1: Селекторы и свойства}

Исходный pen содержит размеченный текст известной микро-сказки
Л.\,Петрушевской <<Пуськи Бятые>>: \textit{<<Сяпала Калуша с
Калушатами по напушке~...>>}. Разметку трогать запрещено --- задача
исключительно в написании четырёх CSS-правил.

\subsection{Задача 1 --- синий заголовок <<Калуша с Калушатами>>}

Главная сложность: на странице два \texttt{<h2>} --- один в шапке
секции (<<Разметка к задачам:>>) и один вложенный в
\texttt{<div>} (<<Калуша с Калушатами>>). Простой селектор
\texttt{h2} перекрасил бы оба. Нужен селектор, попадающий только во
вложенный заголовок:

\begin{lstlisting}[language=CSS]
section > div > h2 {
  color: blue;
}
\end{lstlisting}

\subsection{Задача 2 --- коричневые упоминания Калушат}

Все упоминания Калушат в тексте обёрнуты в \texttt{<strong>}.
Поскольку других \texttt{<strong>} в pen-е нет, достаточно
селектора по тегу:

\begin{lstlisting}[language=CSS]
strong {
  color: brown;
}
\end{lstlisting}

\subsection{Задача 3 --- жёлтая жирная Бутявка в цитатах}

Слово <<Бутявка>> (и его словоформы) встречается в тексте многократно
и обёрнуто в \texttt{<span>}. Нужно стилизовать только те
\texttt{<span>}, что находятся внутри цитат \texttt{<blockquote>} ---
без перекраски Бутявок в обычных абзацах:

\begin{lstlisting}[language=CSS]
blockquote span {
  color: yellow;
  font-weight: bold;
}
\end{lstlisting}

\subsection{Задача 4 --- курсивные Калушата в списке}

В нумерованном списке \texttt{<ol>} тоже встречаются Калушата ---
к ним нужно применить курсивное начертание дополнительно к уже
заданному коричневому цвету:

\begin{lstlisting}[language=CSS]
ol strong {
  font-style: italic;
}
\end{lstlisting}

% — Скриншот: рендер решения блока 2.1 —
\begin{figure}[h]
  \centering
  \ReportImage[width=0.85\textwidth]{img/03_block_2_1}
  \caption{Рендер решения блока 2.1: разноцветный текст, жёлтые Бутявки в цитатах, курсив Калушат в списке}
  \label{fig:block_2_1}
\end{figure}

\begin{figure}[h]
  \centering
  \ReportImage[width=0.85\textwidth]{img/07_codepen_fork_2_1}
  \caption{Форк задания по CSS-селекторам в Codepen}
  \label{fig:codepen_fork_21}
\end{figure}

Готовый форк pen-а на Codepen:
\begin{lstlisting}
https://codepen.io/blxckweed/pen/vEymBmY
\end{lstlisting}

\section{Блок 2.2: Оформление текстовых блоков}

Исходный pen содержит размеченную энциклопедическую статью про
А.~Эйнштейна (заголовок, портрет, лид-абзац, несколько обычных
абзацев, список научных работ, цитата). Задача --- написать пять
CSS-правил.

\subsection{Задача 1 --- главный заголовок}

Заголовок с классом \texttt{main-header}: размер 56\,px, цвет
\texttt{\#365f91} (тёмно-синий), шрифт с засечками.

\begin{lstlisting}[language=CSS]
.main-header {
  font-size: 56px;
  color: #365f91;
  font-family: serif;
}
\end{lstlisting}

\subsection{Задача 2 --- размер всех абзацев}

\begin{lstlisting}[language=CSS]
p {
  font-size: 20px;
}
\end{lstlisting}

\subsection{Задача 3 --- подчёркивание в лид-абзаце}

Внутри абзаца с классом \texttt{lead} (первый абзац --- расширенное
определение) слова \textit{<<Альбе́рт Эйнште́йн>>} выделены тегом
\texttt{<strong>}. Нужно добавить им подчёркивание:

\begin{lstlisting}[language=CSS]
p.lead strong {
  text-decoration: underline;
}
\end{lstlisting}

\subsection{Задача 4 --- два уровня маркеров списка}

Список научных работ Эйнштейна --- \texttt{<ul~class="works-list">}.
Внутри первого \texttt{<li>} есть вложенный \texttt{<ul>} с одним
пунктом про закон взаимосвязи массы и энергии. Требуется:
\begin{itemize}
  \item все пункты обоих уровней --- цвета \texttt{\#365f91};
  \item первый уровень --- маркер <<квадрат>> ($\blacksquare$);
  \item второй уровень --- маркер <<пустой круг>> ($\bigcirc$).
\end{itemize}

\begin{lstlisting}[language=CSS]
.works-list {
  color: #365f91;
  list-style-type: square;
}

.works-list ul {
  list-style-type: circle;
}
\end{lstlisting}

\subsection{Задача 5 --- цитата с фоновой картинкой}

Цитата \texttt{<blockquote>} в конце статьи --- про вклад
Эйнштейна в популяризацию физики. Требовалось:
\begin{itemize}
  \item установить фоновую картинку (по ссылке методички ---
        изображение сетки/паттерна);
  \item изменить цвет текста на \texttt{\#494429};
  \item включить курсивное начертание и шрифт с засечками.
\end{itemize}

Картинка фона лежит локально в \texttt{codepen-solutions/img/blockquote\_bg.jpg}
(скачана с Яндекс.Диска по ссылке методички); в Codepen её можно
загрузить как \texttt{Asset} через
\texttt{Assets~\(\rightarrow\)~Image} и подставить полученный URL.

\begin{lstlisting}[language=CSS]
blockquote {
  background-image: url("../img/blockquote_bg.jpg");
  color: #494429;
  font-style: italic;
  font-family: serif;
}
\end{lstlisting}

% — Скриншот: рендер решения блока 2.2 —
\begin{figure}[h]
  \centering
  \ReportImage[width=0.85\textwidth]{img/04_block_2_2}
  \caption{Рендер решения блока 2.2: крупный заголовок, подчёркнутый лид, цветной список с двумя уровнями маркеров, цитата на фоновой картинке}
  \label{fig:block_2_2}
\end{figure}

\begin{figure}[h]
  \centering
  \ReportImage[width=0.85\textwidth]{img/08_codepen_fork_2_2}
  \caption{Форк задания по оформлению текста в Codepen}
  \label{fig:codepen_fork_22}
\end{figure}

Готовый форк pen-а на Codepen:
\begin{lstlisting}
https://codepen.io/blxckweed/pen/vEymBxB
\end{lstlisting}

\section{Готовые форки в Codepen}

Сводная таблица ссылок на форки, сохранённые под аккаунтом студента:

\begin{center}
\begin{tabular}{|c|l|l|}
\hline
\textbf{Блок} & \textbf{Тема} & \textbf{URL форка} \\
\hline
1.1 & Теги для разметки текста        & \url{https://codepen.io/blxckweed/pen/EaNmYge} \\
1.2 & Списки и таблицы                & \url{https://codepen.io/blxckweed/pen/YPpVKpm} \\
2.1 & Селекторы и свойства            & \url{https://codepen.io/blxckweed/pen/vEymBmY} \\
2.2 & Оформление текстовых блоков     & \url{https://codepen.io/blxckweed/pen/vEymBxB} \\
\hline
\end{tabular}
\end{center}

\textbf{Самопроверка:}
\begin{itemize}
  \item все четыре стартовых pen-а форкнуты под собственным
        аккаунтом Codepen, изменения сохранены;
  \item четыре URL форков подставлены в таблицу выше и в листинги
        внутри глав;
  \item локальные \texttt{index.html} в \texttt{lab12/codepen-solutions/}
        открываются в браузере и визуально совпадают с
        Codepen-форками;
  \item все четыре скриншота (\texttt{01\_block\_1\_1.png} ---
        \texttt{04\_block\_2\_2.png}) положены в
        \texttt{latex-report/img/}, имена совпадают с
        \texttt{\textbackslash includegraphics} в этой главе.
\end{itemize}

\endinput

\input{parts/conclusion}

\end{document}
